\documentclass{article}
\usepackage{amsmath,amssymb,hyperref,verbatim}
\title{Head-Typed Attention Routing for Training-Free Editing\\ on Streaming Video Generators}
\date{}
\begin{document}
\maketitle

\noindent\textbf{Audience note.} This document is written to be self-contained
for an implementing agent working inside the \texttt{DSL-Lab/StreamEdit}
repository. It restates the StreamEdit baseline in full before stating the
delta. Items marked \texttt{CHECK} must be verified against the source code, not
assumed. Items marked \texttt{UNVALIDATED} are design choices with no evidence
behind them yet.

\section{Motivation and problem statement}

\textbf{Task.} Training-free text-driven video editing on a video DiT. Inputs:
source video $V^{s}$, source prompt $p^{s}$, target prompt $p^{t}$, and
optionally an edited first frame $I^{\mathrm{ref}}$. Output: an edited video
adopting the appearance change specified by $p^{t}$ (and $I^{\mathrm{ref}}$)
while retaining the motion and structure of $V^{s}$. No per-video finetuning.

\textbf{Goal (primary, measurable).} On FiVE-Bench: raise edit success
(FiVE-Acc, edit-region CLIP) at equal-or-better background/structure
preservation and \emph{unchanged NFE and s/frame} relative to StreamEdit, and
preserve in-region motion on the non-rigid substitution split --- by replacing
StreamEdit's head-agnostic anchoring with head-typed routing. Every experiment
on the board serves this sentence; anything that does not is out of scope.

\textbf{Gap.} Appearance/motion decoupling is the central mechanism of video
editing. On UNet backbones it is free: spatial self-attention (SA-S) and
temporal self-attention (SA-T) are separate blocks, so one can inject
reconstruction features through SA-S and motion features through SA-T (UniEdit).
Video DiTs replace this with unified 3D full attention, so the split does not
exist to route through. DiT editing methods therefore route along surrogate
axes: \emph{where} (spatial masks), \emph{when} (timestep-dependent blend
ratios), \emph{which layer} (heuristics inherited from UNets). No \emph{editing} method routes along
\emph{what} the injected signal carries. (Head-typed routing has now appeared
for \emph{motion transfer} --- HALO, see prior work --- but on bidirectional
50-step backbones with inner-loop optimization, and with editing only as a
side demo.)

Two consequences:
\begin{enumerate}
  \item A spatial mask cannot preserve the motion of an object it is
        regenerating: inside the mask, source appearance and source motion are
        discarded together.
  \item Preserving structure requires anchoring to the source \emph{everywhere},
        which drags the edited region's appearance back toward the source.
        This is \textbf{under-editing / over-anchoring}, the dominant failure
        mode of StreamEdit and its lineage. Existing systems treat it as a
        tradeoff dialled by $r_{t_i}$ and $\rho$. We claim it is separable.
\end{enumerate}

\textbf{Thesis (revised 2026-07-17).} StreamEdit's anchoring does two jobs at
once. Geometric stability is a \emph{correspondence} constraint and lives in
the temporal heads' attention \emph{map}. Appearance anchoring is a
\emph{value} constraint and lives in the spatial heads' $V$. They are
conflated only because the bridge is head-agnostic. That head type \emph{can}
separate motion from structure in a video DiT is no longer ours to claim ---
HALO demonstrates it training-free on bidirectional backbones, with 50-step
sampling and inner-loop latent optimization. The surviving open cell, and our
claim: \textbf{the separation survives block-causal few-step distillation and
is exploitable by routing alone} --- no optimization, no training, zero extra
NFE --- \textbf{inside an editing pipeline with dual appearance authorities
(source outside $m$, $I^{\mathrm{ref}}$ inside $m$) and background-preservation
constraints}. At 15/5-step streaming budgets, optimization-based head control
is structurally unavailable; routing is the only instrument left, and whether
it suffices is exactly claim~1.

\textbf{Structural evidence that StreamEdit under-edits.} It contains two
independent compensations for weak edits: cross-attention boosting (multiplying
trigger-word attention by $\omega$ inside the edit region) and source-oriented
guidance. An edit-strength amplifier is only added when edits come out weak.

\section{Method}

\subsection{Notation}

\textbf{Self-Forcing architecture (confirmed against \texttt{causal\_model.py}).}
A \textbf{block} is the atomic generation unit: 3 latent frames, attended
\emph{bidirectionally} among themselves. A \textbf{window} is one
\texttt{inference()} call; blocks are generated \emph{sequentially} within it
(block $b$ denoised over its steps, committed to the KV cache at clean $t=0$,
then block $b{+}1$ begins). \texttt{rollout\_inference()} repeats
\texttt{inference()} on overlapping windows for long video.

Crucially, when block $b$ is denoised its query is the current 3 frames and its
keys/values are the \emph{entire visible window} --- all committed frames of
blocks $0,\ldots,b{-}1$ (from the cache) plus block $b$ itself --- attended with
\texttt{causal=False}, i.e.\ \emph{no intra-call mask}. So the current block
attends bidirectionally and directly to every visible past frame, up to
\texttt{local\_attn\_size} frames. The block-causal behaviour comes from
generation order and cache contents, not from a mask inside the attention call.
The asymmetry that ``past cannot attend forward'' concerns earlier blocks'
\emph{queries}, which are frozen and never recomputed; it does not shorten the
current block's attentional reach.

\textbf{Consequence for us:} the stride-$h_l w_l$ temporal stripe is present in
the current block's \emph{actual} attention map, spanning the whole visible
window --- there is no need to construct a separate bidirectional pass to reveal
it. Its \emph{length} equals the number of currently visible frames
$= \min(\text{frames generated in window}, \texttt{local\_attn\_size})$.

\textbf{Cache mechanics (corrected 2026-07-17).} Default
\texttt{local\_attn\_size} $= -1$ $\Rightarrow$ \texttt{max\_attention\_size}
$= 32760$ tokens $= 21$ latent frames of \emph{global} reach; the eviction
branch never triggers. For standard-length runs (FiVE-Bench), one
\texttt{inference()} call covers the whole clip: the cache fills block-by-block
with clean ($t=0$) committed KVs, reach ramps $3 \to 21$ frames monotonically,
and there is no reset --- \textbf{the first cache-saturated block is the last
block of the clip}. The per-window ramp-and-reset model applies only to
long-video \texttt{rollout\_inference()}, whose fresh windows are seeded with
the last overlap block via \texttt{src/trg\_initial\_latent} (overlap frames
cached as clean context, not regenerated). Side effect: the reference frame
conditions \emph{only} window~1; from window~2 the initial-latent slot is
overwritten by the overlap (see risk~5).

Layer $\ell \in [L]$, head $h \in [H]$; $\ell$ suppressed. $S_b = 3 h_l w_l$ =
current block's query tokens. $S_{\mathrm{vis}} =
\min(\text{gen frames},\texttt{local\_attn\_size})\cdot h_l w_l$ = visible key
length. Superscripts $s$, $t$, $\mathrm{ref}$: source, target, reference.
Subscript $p$: cache of prior blocks in the visible window.

$m \in \{0,1\}^{S_{\mathrm{vis}}}$ is the \textbf{edit mask}; $m_j = 1$ iff
token $j$ is in the edit region. (Confirmed: StreamEdit's
\texttt{mask\_bin = mask\_soft > 0}, True $=$ foreground $=$ edit region; the
cached mask is the src$\cup$trg union.)

$r_i = 1 - t_{i-1}^{\rho}$: StreamEdit blend ratio. $t^{*} = 0.5$: injection
threshold. Denoising steps: 15 (standard) / 5 (long-video). $\tau_h$: head
type, \texttt{tau[30,12]} ($L = 30$ layers, $H = 12$ heads).

\subsection{Baseline: what StreamEdit computes}

Dual-branch sampling with shared noise $\epsilon_{t_i}$:
\[
  x^{s}_{t_i} = (1-t_i)\,x^{s}_0 + t_i \epsilon_{t_i}, \qquad
  x^{t}_{t_i} = (1-t_i)\,z^{t}_{t_i} + t_i \epsilon_{t_i}.
\]
The source branch is denoised with $p^{s}$; the target branch with $p^{t}$,
conditioned on the source branch through a self-attention bridge applied
\textbf{identically to every head}:
\[
  \tilde{Q} = r_i Q^{t} + (1-r_i) Q^{s}, \qquad
  \tilde{K} = r_i K^{t} + (1-r_i) K^{s},
\]
\[
  \tilde{K}_p = m_p \odot \big(r_i K^{t}_p + (1-r_i) K^{s}_p\big)
                + \bar m_p \odot K^{t}_p,
\]
\[
  K = \big[\,\tilde{K},\; \tilde{K}_p,\; \mathbb{1}[t<t^{*}]\cdot(\bar m \odot K^{s})\,\big],
  \qquad
  V = \big[\,V^{t},\; V^{t}_p,\; \mathbb{1}[t<t^{*}]\cdot(\bar m \odot V^{s})\,\big],
\]
\[
  O = \mathrm{softmax}\!\big(\tilde{Q}K^{\top}/\sqrt{d}\big)\,V .
\]
The edit mask $m$ comes from cross-attention grounding (thresholded difference
between trigger-word and non-trigger-word attention, unioned over branches).
Cross-attention boosting and source-oriented guidance sit outside the bridge and
are \textbf{left untouched by this proposal}.

Three observations drive the method. $\tilde{Q}$ drags \emph{every} query toward
the source, including inside $m$. $\tilde{K}_p$ drags cache keys toward the
source \emph{inside} $m$. The third block imports source content. The first two
are geometric anchoring bleeding into appearance.

\subsection{Step 1: head classification on the source branch}

\textbf{Where to classify --- block.} Reach ramps $3 \to 21$ frames over the
clip (standard runs; see cache mechanics), so \textbf{classify at the last
block of the clip}, where the temporal stripe spans the full reach. Never
classify at block~0, where a long-range temporal head is indistinguishable from
a spatial one.

\textbf{Sequencing (resolved 2026-07-17).} Routing is needed from block~0 of
the target branch, but the classification signal peaks at the last block ---
after the edit would already have run. The source branch is target-independent
(vanilla \texttt{attention(src\_q, src\_k, src\_v)}), so v1 runs a
\textbf{source-only pre-pass} over the clip to classify, then the dual-branch
edit ($\approx$ half of one extra pass, once per edit). If E1 shows cross-input
stability, replace with per-model offline classification and drop the
pre-pass.

\textbf{Two-axis diagnostic (E1).} The reach ramp confounds a naive flip-rate
measurement, so separate it: (i) \emph{within-window ramp} --- block~0 vs.\ last
block of window~1; head types \emph{will} differ here purely from truncated
reach, which is expected and not a finding; (ii) \emph{across-window stability}
--- last block of window~1 vs.\ last blocks of windows $2,3,\ldots$ (all
saturated); if head type is structural this is near-constant, and drift here is a
genuine signal that forces per-window reclassification. Only axis (ii) tests the
freeze assumption. (Axis (ii) applies to long-video rollout only; standard runs
have a single window.)

\textbf{Where to classify --- denoising step.} The cached past-block keys are
\emph{clean} (committed at $t=0$). The most representative attention map is
clean-query $\times$ clean-key, which occurs only at the \textbf{final
(lowest-noise) step} of the chosen block --- the sole point where the current
block's queries match the clean content the cache holds. High-noise queries are
near-Gaussian and least structured. Classify at the last step; optionally
majority-vote over the last 2--3 steps. (This corrects an earlier draft that
said highest-noise step --- backwards.)

\textbf{No synthetic pass needed.} Because the real forward call already attends
bidirectionally over the full visible window (\texttt{causal=False}), the stripe
is present in the actual attention map. Classify on the source branch's real
attention call; do not build a separate bidirectional reconstruction pass. The
classifying call has the same structure as the calls used during editing, so the
frozen label is valid at apply-time.

\textbf{Freeze.} Classify once at the last block of the first window and freeze
$\tau$ for the whole edit (all blocks, all windows), subject to axis-(ii)
validation above.

\textbf{Criterion.} At the chosen block, sample $\mathcal{P} \subset [S_b]$
current-block query rows, $|\mathcal{P}| = \lceil 0.01\,S_b \rceil$. Define two
candidate key-sets over the visible window: $\mathcal{B}^{\mathrm{spat}}(j)$ =
visible keys in the same latent frame as $j$; $\mathcal{B}^{\mathrm{temp}}(j)$ =
visible keys at the same spatial position $(x,y)$ across all visible frames
(stride $h_l w_l$, up to the visible reach). Wan self-attention contains no
text tokens (text enters via cross-attention), so both sets are purely visual.
Then
\[
  \tau_h = \arg\min_{c \in \{\mathrm{spat},\mathrm{temp}\}}
  \big\|\,\mathrm{Attn}_{\mathcal{B}^{c}}(Q^{s}_{\mathcal{P},h}, K^{s}_h, V^{s}_h)
       - \mathrm{Attn}_{\mathrm{full}}(Q^{s}_{\mathcal{P},h}, K^{s}_h, V^{s}_h)\,\big\|_2^2 ,
\]
where $\mathrm{Attn}_{\mathrm{full}}$ is the real full-visible-window attention.
SVG's output-reconstruction criterion, $\approx 3\%$ overhead, no $O(S^2)$
materialisation. Deliberately \emph{not} FYM's map-shape criterion.

Classification runs on the \textbf{source branch only}; $Q^s, K^s, V^s$ are
already computed in the bridge, so profiling adds no extra forward passes.

\subsection{Step 2: query routing}

\[
  \text{temporal:}\quad Q_h \leftarrow (1-r^{T}_i)\,Q^{s}_h + r^{T}_i\,Q^{t}_h,
  \qquad r^{T}_i := 0 .
\]
\[
  \text{spatial:}\quad \varrho_j = m_j + (1-m_j)\,r_i,
  \qquad Q_{h,j} \leftarrow \varrho_j\,Q^{t}_{h,j} + (1-\varrho_j)\,Q^{s}_{h,j}.
\]
The spatial rule is the under-editing fix: $\varrho_j = 1$ inside the edit
region gives a pure target query with no source anchoring, and outside it
degrades to StreamEdit's schedule, preserving the background.

\subsection{Step 3: key/value routing}

\textbf{Temporal heads} --- correspondence carriers. Keys from source, values
from target, index-aligned:
\[
  K_h = [\,K^{s}_h,\; K^{s}_{p,h}\,], \qquad
  V_h = [\,V^{t}_h,\; V^{t}_{p,h}\,]
  \qquad\Longrightarrow\qquad
  O_h = A^{s}_h V^{t}_h .
\]
The source video's attention map applied to target content. No mask, no
source-$KV$ concat block, no appearance crossing. $A^{s}_h$ is never
materialised: this runs through the same fused SDPA kernel, which matters since
$S \approx 32$k makes explicit map transfer infeasible.

\textbf{Spatial heads} --- appearance carriers. Two appearance authorities,
selected by the mask. Write the concatenated key set as
$\mathcal{K} = \mathcal{K}^{t} \cup \mathcal{K}^{t}_{p} \cup \mathcal{K}^{s}
\cup \mathcal{K}^{\mathrm{ref}}$ and define an \emph{additive} bias (not
elementwise multiplication --- see implementation note below):
\[
  B(j, k) =
  \begin{cases}
    0 & k \in \mathcal{K}^{t} \cup \mathcal{K}^{t}_{p} \\
    0 & k \in \mathcal{K}^{s},\;\; m_k = 0,\;\; t < t^{*} \\
    0 & k \in \mathcal{K}^{\mathrm{ref}},\;\; m_j = 1,\;\; t < t^{*} \\
    -\infty & \text{otherwise,}
  \end{cases}
  \qquad
  O_h = \mathrm{softmax}\!\Big(\tfrac{Q_h \mathcal{K}_h^{\top}}{\sqrt d} + B\Big) \mathcal{V}_h .
\]
$B$ has shape $[B, H_{\mathrm{spat}}, S_q, S_k^{\mathrm{full}}]$ where
$S_k^{\mathrm{full}} = S_q + S_{\mathrm{prev}} + S_q + P$. In code it is three
tensor-index assignments (no loops); see pseudocode.

The asymmetry is load-bearing and invisible in the $\odot$ form:
\textbf{source keys are gated by key index $m_k$; reference keys are gated by
query index $m_j$.} Source keys inside the edit region are denied to
\emph{all} queries ($m_k=1 \Rightarrow -\infty$). Reference keys are denied to
background queries ($m_j=0 \Rightarrow -\infty$). Target keys are always
visible.

\textbf{The mask $m$ is 3D.} $m \in \{0,1\}^{F_{\mathrm{vis}} \times h_l \times w_l}$
where $F_{\mathrm{vis}} = \min(\text{gen frames}, \texttt{local\_attn\_size})$,
flattened to $[S_{\mathrm{vis}}]$ token indices. It is \textbf{not} a static 2D
spatial mask
replicated across frames --- the edited object moves. Three strategies in
increasing cost:
\begin{enumerate}
  \item \emph{Dilated 3D mask (v1 default).} Derive from
        $|\mathrm{VAE}(I^{\mathrm{ref}}) - \mathrm{VAE}(I^{s}_1)|$, threshold,
        then dilate $\times 4$ in both spatial and temporal dimensions within
        each chunk. The 4$\times$ temporal compression of Wan's VAE means one
        latent frame $\approx 4$ pixel frames; a generous spatial dilation absorbs
        within-chunk motion. Crude but avoids external dependencies.
  \item \emph{Flow-warped mask.} Warp the frame-1 mask to each subsequent frame
        using optical flow (RAFT or similar). Offline, per-edit, adds $\sim 2$s.
        Use if v1 degrades on motion-heavy cases.
  \item \emph{Per-chunk re-derived mask.} Diff source-branch latents against the
        reference. Avoids a flow model but introduces drift after chunk~1.
        Avoid unless both above fail.
\end{enumerate}

\textbf{Why spatial heads are the type-correct consumer of $I^{\mathrm{ref}}$.}
A single reference frame has no temporal extent. A temporal head attends to the
same spatial position \emph{across} frames; a one-frame bank gives it nothing to
mix. A spatial head attends \emph{within} a frame, and $I^{\mathrm{ref}}$ is a
frame. The routing is forced by the shape of the operation, not chosen. This
predicts (claim~4) that reference injection into temporal heads actively
degrades results.

\subsection{Step 4: the reference KV bank}

One extra DiT forward on $\mathrm{VAE}(I^{\mathrm{ref}})$ yields
$K^{\mathrm{ref}}, V^{\mathrm{ref}}$ per layer per head, $h_l w_l$ tokens,
cached once.

Confirmed (2026-07-17): computed \textbf{once}, as a clean context forward at
\texttt{context\_noise} $= 0$ --- the same convention as committed cache
blocks. Cache at $t = 0$, not $t^{*}$. Bonus finding from code verification:
current code drops the reference frame after rollout window~1, so the
persistent bank is also the \emph{structural fix for long-video edit fading}
(risk~5) --- a contribution-adjacent property worth stating.

\texttt{UNVALIDATED} --- \textbf{RoPE for reference keys.} Assigning them frame
index $0$ lets 3D RoPE's temporal decay suppress their influence in later frames.
Default v1: replicate the reference keys $F_{\mathrm{vis}}$ times, giving each
copy the querying frame's own temporal index (so the relative temporal rotation
is zero). Cost: $F_{\mathrm{vis}} \cdot h_l w_l$ extra keys, bounded by
\texttt{local\_attn\_size}. Ablate against frame-$0$ indexing.
Note: setting the key's temporal rotation to identity is \emph{not} equivalent
--- 3D RoPE contributes a \emph{relative} rotation $R(\theta_{f'} - \theta_f)$,
so a zero key-rotation still produces a growing relative angle with the querying
frame. True temporal neutrality requires $\theta_{f'} = \theta_f$, i.e.\
replication.

\subsection{Step 5: pseudocode}

\textbf{Implementation notes.}

\textbf{(a) Never implement key masking as $\odot$.} An excluded key multiplied
by zero scores $q\cdot 0 = 0$, and $\exp(0)=1$: excluded keys absorb real
softmax mass and act as attention sinks. Always use $-\infty$ additive bias
passed as \texttt{attn\_mask} to \texttt{F.scaled\_dot\_product\_attention}.

\textbf{(b) Head routing is two separate SDPA calls.} Softmax is per-head;
heads never share denominator or interact inside the attention operator.
Temporal and spatial heads run separate \texttt{sdpa} calls with different key
lengths and then their outputs are scattered back into a $[B,H,S,D]$ tensor:
\begin{verbatim}
O = torch.empty_like(Qt)  # [B,H,S,D]
O[:, T]  = O_temporal     # temporal heads to their index positions
O[:, ~T] = O_spatial      # spatial heads to theirs
\end{verbatim}
No interpolation, no mixing. The downstream $W_o$ projection sees an ordinary
$[B,H,S,D]$ output. This is mathematically identical to running all heads in one
call if the key sets were equal; they are not, which is the only difference.

\textbf{(c) OOD key-length risk.} Wan was trained with temporal and spatial
heads seeing the same key set. Temporal heads now see $\approx 2S$ keys;
spatial heads $\approx 3S+P$. Softmax normalises by row sum regardless of
length, so there is no obvious numerical instability. The risk is
\emph{attention entropy collapse}: a head collapsing onto a single key (near
one-hot softmax), producing pixel artifacts. Monitor per-head entropy in first
runs. If collapse occurs, mild temperature adjustment ($\sqrt{d'} > \sqrt{d}$
for spatial, i.e.\ divide scores by a slightly larger constant) resolves it
without changing the method. The deeper OOD risk is for temporal heads
specifically: $K$ and $V$ come from different branches ($K^s$, $V^t$), a
combination Wan never saw. L2 ablation is the test.

\textbf{(d) Cross-attention boosting is kept and made diagnostic.} The boosting
factor $\omega$ (multiplying trigger-word cross-attention scores inside $m$)
compensates for under-editing. If head routing fixes under-editing structurally,
the optimal $\omega$ should fall toward $1$. Sweep $\omega \in \{1, 1.5, 2, 3\}$
and report the optimum. $\omega\gg1$ still required $\Rightarrow$ under-editing
persists $\Rightarrow$ spatial-head routing is insufficient.

\textbf{(e) Concrete bias construction for spatial heads.}
\begin{verbatim}
# After concatenation along key dim:
# [K_tgt | K_tgt_prev | K_src | K_ref]
#   S_q      S_prev      S_q     P
s0, s1 = 0,    S_q
s1, s2 = s1,   s1 + S_prev
s2, s3 = s2,   s2 + S_q       # K_src block
s3, s4 = s3,   s3 + P         # K_ref block

B = torch.zeros(bsz, H_spat, S_q, s4)
if t < t_star:
    # K_src: deny to edit-region queries (key-index gate)
    B[:, :, :,  s2:s3][:, :, m_flat, :] = -torch.inf
    # K_ref: deny to background queries (query-index gate)
    B[:, :, :,  s3:s4][:, :, ~m_flat,:] = -torch.inf
else:
    B[:, :, :, s2:s4] = -torch.inf   # both blocks off
\end{verbatim}
\texttt{m\_flat} is the 3D mask
$m \in \{0,1\}^{F_{\mathrm{vis}} \times h_l \times w_l}$
flattened to $[S_q]$. Not a static 2D mask; see mask strategy in Step~3.

{\small
\begin{verbatim}
# ---- once per edit ----
# ---- once per edit ----
tau         = classify_heads(src_branch,   # [L,H] bool, True = temporal
                block=first_saturated,     # or last block if clip short
                step=final)                # lowest-noise: clean q x clean cache
K_ref,V_ref = encode_reference(I_ref)      # [L,H,P,D], cached
m_flat      = compute_3d_mask()            # [S_vis] bool, True = edit region

# ---- per block, over its denoising steps (N per config; CHECK 5 vs 15) ----
for i in range(N):
    Qs,Ks,Vs = run_source_branch(x_src, t_i)   # cache QKV per layer
    run_target_branch(x_tgt, t_i, patched_attn)

def patched_attn(l, Qt,Kt,Vt, Qs,Ks,Vs, prev_t, prev_s, t):
    T = tau[l]                         # [H] bool, True = temporal
    r = 1.0 - t_prev ** rho            # StreamEdit schedule

    # --- queries [B,H,S,D] ---
    Q = torch.empty_like(Qt)
    Q[:,  T] = Qs[:, T]                # temporal: full source query (r_T=0)
    rho_q = torch.where(m_flat, 1.0, r).view(1,1,-1,1)
    Q[:, ~T] = rho_q*Qt[:,~T] + (1-rho_q)*Qs[:,~T]  # spatial: free inside m

    # --- temporal heads: source pattern, target content ---
    Kh = cat([Ks[:,T], prev_s.K[:,T]], dim=2)  # source keys
    Vh = cat([Vt[:,T], prev_t.V[:,T]], dim=2)  # target values
    O_T = sdpa(Q[:,T], Kh, Vh, attn_mask=causal_bias())

    # --- spatial heads: target self + src(bg) + ref(edit) ---
    Kh = cat([Kt[:,~T], prev_t.K[:,~T], Ks[:,~T], K_ref[l][:,~T]], dim=2)
    Vh = cat([Vt[:,~T], prev_t.V[:,~T], Vs[:,~T], V_ref[l][:,~T]], dim=2)
    B  = build_spatial_bias(m_flat, S_q, S_prev, P, t)   # see note (e)
    O_S = sdpa(Q[:,~T], Kh, Vh, attn_mask=B)

    O = torch.empty_like(Qt)           # scatter back
    O[:,  T] = O_T
    O[:, ~T] = O_S
    return O
\end{verbatim}
}

\subsection{Ablation ladder}

Each rung adds one change. L0 must reproduce StreamEdit numerically --- that is
the correctness test for the patch.

\begin{center}
\begin{tabular}{llllll}
\hline
 & $r^{T}$ & $\varrho$ & source $KV$ & ref bank & \\
\hline
L0 & $r_i$ & $r_i$ & all heads & none & = StreamEdit \\
L1 & $0$ & $r_i$ & all heads & none & + head-typed $Q$ \\
L2 & $0$ & $r_i$ & spatial, $\bar m$ & none & + head-typed $KV$ \\
L3 & $0$ & $m + \bar m r_i$ & spatial, $\bar m$ & spatial, $m$ & \textbf{full} \\
L* & \multicolumn{4}{l}{L3 with \texttt{tau} randomly permuted, size-matched} & control \\
\hline
\end{tabular}
\end{center}

Variant A $=$ L2 (head routing only, $I^{\mathrm{ref}}$ left in StreamEdit's
prev-chunk cache). Variant B $=$ L3. Reporting only B makes the gain
unattributable.

\section{Novelty claims}

\begin{enumerate}
  \item \textbf{(Reworded 2026-07-17, post-HALO.) The head-typed separation of
        geometric and appearance anchoring survives block-causal few-step
        distillation and is exploitable by routing alone --- no inner-loop
        optimization, no training, zero extra NFE.} HALO establishes head-typed
        motion/structure control on bidirectional 50-step backbones \emph{with}
        latent optimization; at 15/5-step streaming budgets that instrument is
        structurally unavailable, and no head-aware method does editing under
        preservation constraints. Operationally: freeing spatial heads inside
        $m$ while fully anchoring temporal heads yields stronger edits at equal
        or better structural stability than StreamEdit, at unchanged NFE and
        s/frame. \emph{Falsified} if structure degrades once spatial heads are
        freed even with temporal heads fully anchored --- geometry not cleanly
        localised in temporal heads under this distillation. \textbf{Primary
        claim and primary falsifier.} Tested by L2/L3 vs L0, supported by a
        cost table against optimization-based head control (HALO-style) and
        trained streaming editing (LiveEdit).
  \item \textbf{Head type is a usable routing variable, not a proxy for
        injecting less.} \emph{Falsified} if L* (size-matched random head
        permutation) matches L3.
  \item \textbf{Head-conditional injection preserves in-region motion for
        non-rigid object substitution, which mask-conditional injection
        structurally cannot.} Measured as tracklet-similarity motion fidelity
        restricted to $m$, on FiVE-Bench's non-rigid substitution split.
  \item \textbf{A single reference frame can only be read by spatial heads.}
        Injecting $K^{\mathrm{ref}},V^{\mathrm{ref}}$ into temporal heads should
        \emph{degrade} results, not merely fail to help. \emph{Falsified} if it
        is neutral or beneficial, which would indicate the head taxonomy does
        not correspond to the operations it is named for.
  \item \textbf{Head type is stable enough across steps, prompts, and branches
        to be a fixed per-input assignment.}
        \texttt{SCREENING-VALIDATED (R9, 2026-07-18)}: bimodal margins; 10-head
        core ($m > 0.4$) with step-flip $0$ and video-flip $4/100$ head-pairs
        (95\% CI $\approx [1\%, 10\%]$, $n = 5$ videos --- point-estimate pass,
        not certification). Frozen per-model \texttt{tau} adopted; no noise
        gating; per-edit pre-pass retained as fallback. Full validation
        accumulates ride-along (criterion $\approx 3\%$ overhead on the source
        forward) across R10/R12 runs before this claim is printed.
\end{enumerate}

\section{Relation to prior work}

\textbf{HALO} (2607.11081, ECCV 2026; posted 2026-07-13, verified 2026-07-17).
\emph{Closest prior art; displaces GenHOI.} Training-free head-typed control
in video DiTs: motion heads identified via optical-flow-aligned displacement
maps, structure heads via attention entropy ($\tau < 7$); reference $V$
injected through structure heads, motion imposed by \emph{latent optimization}
on motion-head cues (50 steps, optimization in 12). CogVideoX $+$ Wan,
bidirectional. Includes an editing \emph{demo} (subject swap, layout kept) ---
no masks, no preservation metrics, no reference frame, no benchmark. Differs
from us: optimization-based, bidirectional multi-step, generation-first. It
hands us the taxonomy-existence prior on Wan (de-risks E0) and takes ``first
head-typed training-free routing in a video DiT'' off the table --- claim~1 is
now regime-specific (see 2026-07-17 changelog).

\textbf{LiveEdit} (ECCV 2026; repo verified 2026-07-17). Streaming
chunk-causal editing built on Wan2.1/Self-Forcing --- \emph{trained} (three
stages: foundation tuning, teacher-forced causal adaptation, DMD compression).
The streaming-editing niche now has a tuned occupant; StreamEdit and this
proposal remain the training-free side. Cite as the trained baseline.

\textbf{MotionAdapter} (2601.01955). Attention-derived motion fields
rearranged via DINO semantic correspondence to guide DiT denoising. Not
head-typed; relevant as the candidate mechanism for softening the
positional-correspondence assumption of our temporal-head map transfer
(risk~3).

\textbf{UniEdit} (2402.13185). Tuning-free video motion+appearance editing;
injects source features through SA-S and motion through SA-T. \emph{Same routing
idea, on a UNet where SA-S/SA-T are separate blocks.} Our contribution is not the
routing but its \emph{recovery} in a backbone where the split does not exist.
State it that way; the reverse framing is sunk by UniEdit.

\textbf{Sparse VideoGen} (2502.01776, ICML 2025). Source of the spatial/temporal
head taxonomy and of the classification criterion we adopt. Used there for
inference acceleration only. Classifies \emph{online} because patterns are
dynamic across steps and prompts.

\textbf{Follow-Your-Motion} (2506.05207, ICLR 2026). Applies the taxonomy to
\emph{tuned} motion transfer via head-typed LoRAs. We take its channel-surgery
construction as reference and reject its offline count-balanced classifier and
its dual-attention-fusion re-parameterisation. Its ablation omits our claim~2.

\textbf{StreamEdit / StreamGVE} (2605.21466, ECCV 2026). Our base. Routes on
\emph{where} and \emph{when}; no component is head-aware. Our delta is the
\emph{what} axis, plus a dedicated path for $I^{\mathrm{ref}}$.

\textbf{GenHOI} (2603.06048; abstract verified 2026-07-17). Injects reference
tokens into video tokens through self-attention with Head-Sliding RoPE
(head-specific temporal offsets countering 3D RoPE decay) and a spatial
attention gate. Trained, domain-specific (hand-object interaction); its
head-specific offsets are positional bookkeeping, not a functional taxonomy.
Displaced by HALO as closest prior art; still cite for the reference-key RoPE
question it raised (Step~4).

\textbf{Analysis of Attention in VDiTs} (2504.10317). Zero-shot editing by
self-attention map transfer on Wan2.1-T2V-1.3B; identifies layers with distinct
roles. Our idea at \emph{layer} granularity. We must argue head granularity is
the right one.

\textbf{ContextFlow} (2509.17818). Argues UNet layer-selection heuristics are
unsuitable in DiTs and that direct transplant causes geometric and temporal
artifacts. Third-party support for our motivation.

\textbf{LoRA-Edit} (2506.10082). First-frame-guided editing via mask-aware LoRA
on Wan2.1-I2V. Already performs the appearance/motion decomposition --- it trains
on isolated edited frames ``to avoid including motion information'' --- but
routes it with \emph{masks}, in one monolithic LoRA. Relevant if the tuned
variant is revived.

\textbf{FREE-Edit} (2603.01164). Training-free image-driven video editing with
adaptive key replacement; reports cross-attention masks fail to localise object
addition. Motivates our alternative mask derivation.

\section{Open questions and risks}

\begin{enumerate}
  \item \textbf{PRIMARY RISK. Five denoising steps may require the anchoring we
        are removing.} StreamEdit anchors aggressively because a few-step sampler
        has little room to recover. Freeing spatial heads inside $m$ may trade
        under-editing for an unanchored blob: wrong placement, unclosed edges,
        chunk-to-chunk incoherence. The thesis is that temporal-head
        correspondence supplies the geometric scaffold alone. If it does not,
        the method reduces to a worse tradeoff dial. \emph{Test with L2 before
        building the reference bank.}
  \item \textbf{Does the head taxonomy survive block-causal attention?}
        Self-Forcing and LongLive are chunked-KV-cache, block-causal, few-step
        distillations. SVG's temporal-head slash signature at stride $h_l w_l$
        was characterised on \emph{bidirectional dense} attention. It may be
        truncated or absent. \emph{Step 0: dump raw $Q,K$ per head per step for
        a few FiVE-Bench cases and look at the maps before writing a
        classifier.} If deformed, re-deriving the pseudo-masks for the chunked
        layout is itself a contribution. \emph{Update 2026-07-17:} HALO confirms
        head specialization on (bidirectional) Wan, raising the prior that the
        taxonomy exists; what remains untested --- and is now the actual claim
        --- is its survival under causal few-step distillation.
  \item \textbf{Temporal-head map transfer assumes positional, not semantic,
        correspondence.} $A^{s}_h$ encodes ``token at $(x,y,f)$ pulls from
        $(x,y,f\pm 1)$ with these weights.'' Applying it to $V^{t}$ presumes the
        edited object occupies the source object's coordinates. Holds for rigid
        substitution with matched silhouette (bus $\to$ jeep); exactly what fails
        for non-rigid (dancer $\to$ horse), which is where claim~3 promises most.
        \emph{Test rigid first.} A dilated mask on temporal heads may be
        needed; MotionAdapter's DINO-correspondence rearrangement is the
        candidate soft fix if dilation is insufficient.
  \item \textbf{$Q^{s}$ and $Q^{t}$ may not be comparable.} StreamEdit's $r_i$
        ramps from $0$, so its query swap is strongest at high noise where the
        latents are most alike. $r^{T}=0$ throughout removes that safety. Measure
        the constant first; then try a slow ramp.
  \item \textbf{Single-frame reference drifts over long video.} Over 470 frames
        the object rotates and disoccludes; a frame-1 bank cannot say what its
        back looks like (LoRA-Edit \S3.3.2's motivation). Plan a multi-frame
        bank. Do \emph{not} re-anchor on generated chunks --- that compounds
        drift.
  \item \textbf{The previous-chunk KV cache should already be head-typed.}
        Spatial heads attend within-frame and should have little use for the
        previous chunk. If the classifier agrees, head-conditional cache
        treatment is a free extra result and evidence the taxonomy is real. If it
        disagrees, that is evidence for risk~2.
  \item \textbf{(Rollout only.) Temporal reach is bounded per window.} Standard
        runs have global 21-frame reach (see cache mechanics), so this risk is
        confined to long-video rollout: there, source motion with period
        exceeding the window cannot be captured by temporal-head injection
        inside one window, and cross-window coherence rides on the 3-frame
        overlap and the target branch's own continuity. Inherited from
        StreamEdit. Upside: bounded reach caps the spatial-head key length, so
        the two-SDPA cost does not grow with total video length.
  \item \textbf{Recombination objection, now concrete (revised 2026-07-17).}
        The steelman a reviewer will write: ``SVG's taxonomy $+$ UniEdit's
        routing $+$ HALO's training-free head-typed demonstration $+$
        StreamEdit's scaffold --- what did you discover?'' Answerable only by
        results, on three legs: (i) a \emph{capability} neither parent shows
        --- non-rigid in-region motion preservation with strong edits at
        streaming speed (claims 1 and 3); (ii) the L* control showing head
        identity is causal (neither HALO nor FYM ran it); (iii) the cost table
        showing optimization-based head control is orders of magnitude off the
        streaming budget. If L2/L3 merely match StreamEdit on FiVE-Bench, there
        is no paper on this thesis --- with HALO public, the framing fallback
        is gone. Concurrent-work pressure (two adjacent ECCV'26 papers, one
        posted this week) argues for pace.
\end{enumerate}

\subsection*{Code-verification TODO (do before implementing)}

All items resolved 2026-07-17 against the Self-Forcing build; see the
2026-07-17 changelog entry for consequences. File references are to
\texttt{Self-Forcing\_StreamEdit/}.

\begin{enumerate}
  \item \emph{Resolved.} Seeded, not regenerated:
        \texttt{rollout\_inference} feeds the last overlap frames as
        \texttt{src/trg\_initial\_latent} to the next \texttt{inference()} call
        (\texttt{edit\_causal\_inference.py:180--181}), which caches them as
        clean context. Overlap frames are stripped from the output. Default
        overlap $=$ 1 block (3 latent frames). \textbf{Side finding: the
        reference frame is dropped after window~1} --- \texttt{initial\_latent}
        is overwritten by the overlap for windows $\ge 2$. Sharpens risk~5; the
        dedicated ref bank fixes it structurally.
  \item \emph{Resolved.} \texttt{local\_attn\_size} $= -1$ by default
        (\texttt{utils/wan\_wrapper.py:121}) $\Rightarrow$
        \texttt{max\_attention\_size} $= 32760$ tokens $= 21$ latent frames of
        \textbf{global} reach (\texttt{causal\_model.py:127}); the eviction
        branch requires \texttt{local\_attn\_size} $\ne -1$, so standard runs
        have no rolling window and no reset. Steps: 15 (standard) / 5
        (long-video). \texttt{num\_frame\_per\_block} $= 3$
        (\texttt{configs/self\_forcing\_dmd.yaml:48}).
  \item \emph{Resolved.} \texttt{mask\_bin = mask\_soft > 0}
        (\texttt{edit\_causal\_inference.py:755}); \textbf{True $=$ foreground
        $=$ edit region}. The cached \texttt{trg\_fg\_mask} is the union
        $\mathrm{src} \cup \mathrm{trg}$ fg.
  \item \emph{Resolved.} Head-agnostic confirmed: every bridge term
        (\texttt{causal\_model.py:315--375}) indexes batch/token dims only;
        \texttt{blender\_rate} is scalar.
  \item \emph{Resolved.} $V^{s}$ is used: masked background
        \texttt{src\_current\_value} is appended alongside keys
        (\texttt{causal\_model.py:355}).
  \item \emph{Resolved.} Gather $+$ concat and in-place blend on masked rows;
        \textbf{no $\odot$ zeroing anywhere} --- the speculated attention-sink
        bug does not exist in the baseline. Note~(a) applies only to our new
        bias-based implementation.
  \item \emph{Resolved.} Once. Context/VP KV is a single clean forward at
        \texttt{timestep = context\_noise} $= 0$
        (\texttt{configs/default\_config.yaml:10},
        \texttt{edit\_causal\_inference.py:322--376}), same convention as
        committed cache blocks. Cache the ref bank at $t=0$, not $t^{*}$.
  \item \emph{Resolved.} With \texttt{sink\_size} $= 0$ (the live path), the
        cache stores \textbf{unroped} keys; RoPE is re-applied to the whole
        visible cache at every attention call with frames indexed from 0 and the
        query re-based via \texttt{last\_chunk\_start\_frame}
        (\texttt{causal\_model.py:280--295}). The reference bank's frame-index
        assignment (replication scheme and frame-$0$ arm) is therefore rope-time
        bookkeeping, not a cache-format change.
  \item \emph{Resolved.} $H = 12$, $L = 30$
        (\texttt{wan/configs/wan\_t2v\_1\_3B.py}); \texttt{tau} is $[30, 12]$
        --- 360 heads total, small enough for greedy per-head oracle search as
        an analysis experiment.
\end{enumerate}

\section{Changelog}

\subsection*{2026-07-10 --- Initial document; anchored on StreamEdit}

Seeded from Follow-Your-Motion (2506.05207), LoRA-Edit (2506.10082), StreamEdit
(2605.21466).

Two directions considered. \emph{(a) Tuned}: extend LoRA-Edit with head-typed
LoRAs --- appearance LoRA on spatial heads trained on $I^{\mathrm{ref}}$, motion
LoRA on temporal heads trained on $V^{s}$, both loaded at inference (inverting
Follow-Your-Motion, which discards its spatial LoRA). \emph{(b) Training-free}:
make StreamEdit's bridge head-conditional.

\textbf{(a) set aside.} Its delta is thin: LoRA-Edit already performs the
appearance/motion decomposition (its \S3.3.2 trains on isolated edited frames
``to avoid including motion information'' --- Follow-Your-Motion's single-frame
spatial stage under another name) and routes it with masks. The only difference
(a) offers is head routing instead of mask routing, which is the \emph{same
question} (b) asks, bundled with per-video tuning cost, and with
Follow-Your-Motion already holding the flag on head-typed LoRA in a DiT. Revive
only if (b) establishes head routing works and stronger reference-frame
appearance control is wanted.

\textbf{(b) adopted.} Base: StreamEdit, Self-Forcing variant, Wan2.1-T2V-1.3B.
Classification via SVG's output-reconstruction criterion rather than
Follow-Your-Motion's map-shape criterion --- functional rather than cosmetic
definition; avoids materialising the $O(S^2)$ map; already online and per-step,
which the dual-branch setting needs. Classification on the source branch only.
Evaluation on FiVE-Bench.

\textbf{Dropped:} hand-specifying injection rules per edit type (``inject spatial
for colour change, not for object substitution''). With six edit types and a
binary switch it is a hyperparameter fitted to FiVE-Bench, not a mechanism. If
the decomposition is real, edit-type behaviour should fall out of the
mask~$\times$~head-type~$\times$~timestep product. If it does not, that is
evidence \emph{against} the decomposition and should be reported as such.

\subsection*{2026-07-10 --- Integration mechanism: map-vs-values, not fusion}

\textbf{Follow-Your-Motion's dual attention fusion rejected.} It is an exact
identity re-parameterisation of the attention block whose sole function is to
expose per-head-type linear layers for LoRA. Training-free, it computes the same
function at extra cost. Revive only with direction (a).

\textbf{A flaw in the initial entry was corrected.} It proposed injecting source
$K$ \emph{and} $V$ for temporal heads everywhere including inside $m$. Since $V$
carries content, this would import the source object's appearance into the edited
region --- the exact failure the method exists to avoid. Corrected routing:
the attention map is correspondence, the values are appearance. Temporal heads
take $Q^{s},K^{s},V^{t}$ (implicit head-selective map transfer through the fused
kernel, no $O(S^2)$ materialisation).

Formulated as per-head-type blend coefficients so StreamEdit is strictly nested,
making the size-matched random-head control a one-line permutation.

\subsection*{2026-07-10 --- Reframed around under-editing; reference frame given its own path}

\textbf{Error 1: $I^{\mathrm{ref}}$ was omitted entirely.} It is not a peripheral
input but a second appearance authority. In StreamEdit it enters only as a
previously generated chunk in $K^{t}_p,V^{t}_p$, where $\tilde{K}_p$ blending
dilutes it back toward $K^{s}_p$ inside the mask. It now gets a dedicated path:
spatial heads read $I^{\mathrm{ref}}$ inside $m$, the source video outside $m$.

\textbf{Error 2: spatial heads were left on StreamEdit's recipe.} That recipe ---
query blending toward $Q^{s}$ everywhere, masked source-$KV$ concat --- \emph{is}
the cause of under-editing/over-anchoring. Leaving it intact would have preserved
the exact pathology the method exists to remove. Spatial heads inside $m$ now
receive no source anchoring of any kind.

\textbf{Thesis reframed.} StreamEdit's anchoring does two jobs at once:
geometric stability (correspondence, in temporal heads' attention map) and
appearance anchoring (values, in spatial heads' $V$), conflated only because the
bridge is head-agnostic. The contribution is that head type separates them.
Separability is now the primary novelty claim and the primary falsifier,
demoting the earlier ``orthogonal routing variable'' framing.

\textbf{New: type-correctness argument.} A single reference frame has no temporal
extent, so only within-frame heads can read it. The routing of $I^{\mathrm{ref}}$
to spatial heads is forced by the shape of the operation. Yields the falsifiable
prediction that reference injection into temporal heads actively degrades.

\textbf{New prior art: GenHOI} (2603.06048), head-wise treatment of
reference-token self-attention injection (Head-Sliding RoPE). Trained and
domain-specific, so it does not scoop the proposal, but it is the closest
mechanism-level precedent and it raises the RoPE question for reference keys,
previously unconsidered.

\subsection*{2026-07-10 --- Restructured for implementation handoff}

Document reorganised for an agent working in the StreamEdit repo. No decision
was changed. \S2 now runs notation $\to$ baseline $\to$ classification $\to$
query routing $\to$ KV routing $\to$ reference bank $\to$ pseudocode $\to$
ablation ladder; the StreamEdit baseline is restated in full so the file stands
alone. Prior-work entries reduced to one-liners stating the \emph{difference}.
Earlier changelog prose tightened; no recorded decision or reason removed.

Added: the L0--L3/L* ablation ladder, with L0 required to reproduce StreamEdit
numerically as the patch correctness test. Added a Code-verification TODO list.
Added two implementation warnings: (i) key masking must be an additive $-\infty$
bias or a gather, never $\odot$ --- zeroed keys score $0$, so $\exp(0)=1$ and
they absorb softmax mass as degenerate sinks, which may be a latent bug in the
baseline itself; (ii) head routing is an index over an existing axis, requiring
no architectural change. Added the frame-aligned reference-key replication scheme
as the v1 answer to RoPE decay, to be ablated against frame-$0$ indexing.

\subsection*{2026-07-10 --- Spatial-head gating disambiguated; RoPE default corrected}

Two fixes to the implementation spec, no change to the thesis.

\textbf{Gating asymmetry made explicit.} The previous $K^{s}|_{\bar m}$,
$K^{\mathrm{ref}}|_{m}$ restriction notation was ambiguous: it read as if both
blocks were subset by key index, but reference keys carry no mask index. Replaced
by an explicit additive bias $B(j,k)$. Source keys are denied by \emph{key} index
($m_k = 1 \Rightarrow -\infty$); reference keys are denied by \emph{query} index
($m_j = 0 \Rightarrow -\infty$). Both directions are required: the first stops
source appearance re-entering the edit region, the second stops the edited object
leaking into the background. Only the first was previously stated.

\textbf{RoPE for reference keys.} A tempting shortcut --- set the reference keys'
temporal rotation to the identity so they appear ``frame-neutral'' --- is
\emph{wrong}. 3D RoPE contributes $q^{\top} R(\theta_{f'} - \theta_{f})\,k$, a
\emph{relative} rotation. An identity rotation on the key leaves the relative
angle equal to $-\theta_f$, which still grows with the querying frame index, so
the decay is not removed, merely re-signed. True temporal neutrality requires
$\theta_{f'} = \theta_{f}$, i.e.\ replicating the reference keys per querying
frame. That is the v1 default in Step 4 and it stands; the identity-rotation
variant is not a valid ablation arm and has been struck.

\subsection*{2026-07-10 --- Step-by-step clarifications: classification schedule, 3D mask, bias, two-SDPA}

Fifth iteration. No change to the method; four implementation ambiguities
resolved.

\textbf{Classification schedule.} Classify once at $t_N$ of the first window,
at the last block of that window (see sixth iteration below for the correction).
Freeze for the entire edit.

\textbf{3D mask.} $m \in \{0,1\}^{S_w}$ over the window's token sequence
(not a 2D spatial mask replicated across frames). v1: dilated VAE-diff mask.
v2: flow-warped mask. v3: per-block re-derived (avoid). Flattened to $[S_w]$
token indices for the bias construction.

\textbf{Bias construction.} Concrete tensor-index assignments added in note~(e)
and pseudocode. Key asymmetry: source-block denial is a \emph{key}-index gate
($m_k$); reference-block denial is a \emph{query}-index gate ($m_j$). Two
index expressions, no loops.

\textbf{Two-SDPA design.} \texttt{scatter\_heads} is two index assignments into
an empty $[B,H,S,D]$ tensor. No softmax interaction between head groups; the
output is in-distribution for $W_o$. The real OOD risk is in temporal heads
seeing a cross-branch $(K^s, V^t)$ pairing Wan never trained on; L2 is the test.

\textbf{Cross-attention boosting} kept and made diagnostic: if optimal $\omega$
falls toward 1 after head routing, that is supporting evidence for claim~1.

\subsection*{2026-07-10 --- Block/window/rollout terminology fixed; Step~1 classification point corrected}

Sixth iteration. ``Chunk'' was used inconsistently for both the 3-frame block
and the multi-block window. Now standardised: \textbf{block} = 3 latent frames,
bidirectional within; \textbf{window} = $B_w$ causally-ordered blocks in one
\texttt{inference()} call; \textbf{rollout} = repeated overlapping windows.

\textbf{Classification point corrected.} The fifth-iteration entry said
``full window, bidirectionally.'' This was wrong. Within a window, blocks are
causally ordered via KV cache: block $b$'s queries see blocks $0,\ldots,b$ as
keys but not the other way. The temporal stripe pattern (stride $h_l w_l$) lives
in the accumulated KV cache, not in the local $QK$ square of a single block.
Correct recommendation: classify at the \emph{last block} of the first window,
where the cache spans all prior frames and the temporal pattern is maximally
expressed. The E1 diagnostic (flip rate across blocks) is now concrete:
$<5\%$ flip rate $\Rightarrow$ early classification suffices.

Open question marked \texttt{CHECK}: whether Self-Forcing's \texttt{inference()}
uses a sequential growing-cache architecture (assumed above) or a full
within-window bidirectional mask. The classification recommendation changes
depending on the answer.

\subsection*{2026-07-10 --- Attentional reach corrected: full visible window, not 3-frame block}

Seventh iteration. The sixth-iteration model of within-window attention was
wrong and is corrected against \texttt{causal\_model.py}.

\textbf{Correction.} Blocks are generated sequentially, but when a block is
denoised it attends with \texttt{causal=False} over the \emph{entire visible
window} (own 3 frames + all committed past frames, up to
\texttt{local\_attn\_size}), with no intra-call mask. The sixth-iteration claim
that a block sees past frames ``only as cache keys, pattern truncated to 3
frames'' was false: the current block's attentional reach is the full visible
window. The ``past cannot attend forward'' asymmetry concerns frozen earlier
\emph{queries} and does not shorten the current block's stripe.

\textbf{Three consequences.}
(1) \emph{No synthetic bidirectional pass is needed.} The temporal stripe is
present in the real forward attention map. Classify on the actual call. Removes
a chunk of planned complexity.
(2) \emph{Classification block $=$ first cache-saturated block} (visible frames
$= \texttt{local\_attn\_size}$), not ``last block of first window.'' Reach grows
until saturation, then holds; block~0 gives a 3-frame reach that misclassifies
long-range temporal heads as spatial. Short clips: last block.
(3) \emph{Classification step $=$ final (lowest-noise) step}, not highest-noise
as the fifth/sixth iterations stated. Cached keys are clean; only the last step
gives clean-query $\times$ clean-key. This reverses the earlier instruction.

Notation updated: $S_c/F_c$ (chunk) replaced by $S_{\mathrm{vis}}/F_{\mathrm{vis}}$
(visible window, bounded by \texttt{local\_attn\_size}). Per-block step count
flagged \texttt{CHECK} ($5$ assumed earlier vs.\ $15$ read from code).

\subsection*{2026-07-10 --- Cache mechanics: per-window, ramps to \texttt{local\_attn\_size}, resets}

Eighth iteration. Clarified that the window spans exactly
\texttt{local\_attn\_size} frames; the KV cache has matching capacity, is
allocated per window, fills block-by-block with clean committed KVs, and resets
each (overlapping) window. No sliding within a window --- fill reaches capacity
at the last block.

\textbf{Reconciles iterations 6 and 7.} The classification location ``last block
of the first window'' (iteration~6) and ``first cache-saturated block''
(iteration~7) are the \emph{same} block here, because window length $=$
\texttt{local\_attn\_size}. Iteration~6 had the right location for the wrong
reason (it wrongly believed local attention could not see the cross-frame
pattern); iteration~7's saturation argument is the correct justification. No open
contradiction.

\textbf{Added.} (i) Bounded-reach limitation: temporal receptive field $\le
\texttt{local\_attn\_size}$, ramps per window, resets per window; source motion
with longer period is uncapturable within one window. Upside: bounds spatial-head
key length, so two-SDPA cost is independent of total video length. (ii) E1 split
into within-window ramp (expected drift, not a finding) and across-window
stability (the real freeze test). (iii) \texttt{CHECK} on overlap-seeding vs.\
regeneration at window reset, and on the \texttt{local\_attn\_size} value ---
a small value would weaken the taxonomy itself.

\subsection*{2026-07-17 --- Code verification closed (T0/R8); cache model corrected for standard runs}

Ninth iteration. All nine Code-verification TODO items resolved against
\texttt{Self-Forcing\_StreamEdit}; answers recorded inline in the TODO list.
Method unchanged; four consequences for the spec:

\textbf{(1) The eighth-iteration cache model only describes the rollout path.}
Default \texttt{local\_attn\_size} $= -1$ maps to a 21-latent-frame
\emph{global} attention budget with no eviction and no per-window reset.
For standard-length runs (FiVE-Bench), one \texttt{inference()} call covers the
whole clip: reach ramps $3 \to 21$ frames monotonically and \textbf{the first
cache-saturated block is the last block of the clip}. Risk~7's ``small
\texttt{local\_attn\_size} weakens the taxonomy'' is moot for standard runs and
relevant only to long-video rollout. E1's across-window axis applies to rollout
only.

\textbf{(2) Classification sequencing gap made explicit, and its resolution.}
Routing is needed from block~0 of the target branch, but the classification
signal is fully expressed only at the last block, final step --- after the edit
would already have run. Resolution: the source branch is target-independent
(vanilla \texttt{attention(src\_q, src\_k, src\_v)},
\texttt{causal\_model.py:332}), so a \textbf{source-only pre-pass} over the clip
can classify before the dual-branch edit ($\approx$ half of one extra pass, once
per edit). The cheaper alternative --- per-model offline classification frozen
across inputs --- is admissible iff E1 shows cross-input stability; keep the
pre-pass as fallback.

\textbf{(3) Baseline exonerated on the $\odot$ question.} Masking is
gather$+$concat and in-place blends; no zeroed keys, no attention-sink bug to
report. Note~(a) remains binding for our bias-based reimplementation only.

\textbf{(4) Risk 5 is worse than stated, and the ref bank is its fix.} In
current code the reference frame conditions \emph{only window 1}; from window 2
onward \texttt{initial\_latent} is overwritten by the rollout overlap, so
long-video appearance persistence rides entirely on 3 overlap frames. The
dedicated persistent ref bank (Step~4) is therefore not just a routing
convenience but the structural fix for long-video edit fading --- worth stating
as a contribution-adjacent property.

Also recorded: cache keys are stored unroped and re-roped per call (so the
replication RoPE scheme is rope-time bookkeeping); ref/context KV is cached once
at \texttt{context\_noise} $= 0$ (cache the ref bank at $t = 0$, not
$t^{*}$ --- corrects Step~4's fallback suggestion); $H{\times}L = 12 \times 30 =
360$ heads, small enough for a greedy per-head oracle-routing search as an
analysis experiment (upper-bounds any classifier).

\subsection*{2026-07-17 --- Literature refresh: HALO forces repositioning; goal restated; body consistency pass}

Tenth iteration. A litsearch pass (run under a novelty check) found three works
the 2026-07-10 scan missed:

\textbf{HALO} (2607.11081, ECCV'26, posted 2026-07-13 --- four days before
this entry): training-free head-typed motion/structure control in video DiTs
--- motion heads via flow-aligned displacement, structure heads via attention
entropy, reference-$V$ injection through structure heads, latent optimization
on motion-head cues; CogVideoX $+$ Wan, bidirectional, 50 steps; editing shown
as a demo. \textbf{LiveEdit} (ECCV'26): \emph{trained} streaming editing on
Self-Forcing (three-stage distillation) --- the streaming-editing niche now
has a tuned occupant. \textbf{MotionAdapter} (2601.01955):
DINO-correspondence motion-field rearrangement --- candidate risk-3
mitigation. GenHOI verified and displaced as closest prior art.

\textbf{Repositioning.} ``Head type separates motion from structure,
training-free, in a video DiT'' is no longer claimable --- HALO holds it.
Claim~1 reworded to the surviving cell: the separation \emph{survives
block-causal few-step distillation} and is exploitable by \emph{routing alone}
(no optimization, zero extra NFE) in an \emph{editing} pipeline under
preservation constraints. ``Optimization-free / few-step-compatible'' promoted
from implementation detail to headline; title extended to name the streaming
regime; an explicit measurable Goal added to \S1. Risk~8 rewritten as the
concrete HALO-recombination steelman (answer: capability $+$ L* control $+$
cost table; no framing fallback remains). Upside recorded: HALO confirms head
specialization on Wan, de-risking E0. Claim~1's evidence plan gains a cost
table vs.\ optimization-based head control and trained streaming editing.

\textbf{Body consistency pass} (declared restructuring, no method change):
T0-resolved answers folded inline --- cache mechanics rewritten for standard
runs (global 21-frame reach; ramp-and-reset applies to rollout only);
classification point $=$ last block of clip; the block-0 sequencing gap and its
source-only pre-pass resolution added to Step~1; mask-sign, step-count,
ref-KV-once, and no-text-tokens corrections applied; risk~7 scoped to rollout.
Prior-work section reordered to lead with HALO.

\subsection*{2026-07-18 --- R9/E0 executed: taxonomy survives distillation; GO at screening level}

Eleventh iteration. E0 (workboard R9) ran on 5 FiVE clips, all blocks $\times$
steps $\{0,7,12,13,14\}$, SVG output-reconstruction criterion, source branch of
a degenerate edit. \textbf{Verdict: GO.}

\textbf{Results.} Margins bimodal: 281/360 heads strongly spatial
($m < -0.2$), 26 temporal-side, thin middle ($\sim$53). Temporal heads
concentrate at layers 0--9 and 28--29; layers 10--27 essentially empty.
Mechanistic probe confirms semantics: L29H09 is a near-pure
\emph{copy-previous-frame} head (100\% mass on same position, one frame back,
one-hot); L01H04 a soft temporal kernel ($\sim$0.9 stripe mass over recent
frames); the ambiguous band is genuinely spatiotemporal-local. Stability
decomposes by margin strength: $m > 0.4$ core (10 heads) has step-flip 0 and
video-flip $4/100$ pairs; instability is confined to the weak tail
($m \in [0.15, 0.3]$) and the ambiguous band. Block ramp shows margin
compression, not sign flips.

\textbf{Adopted for Step 2/3 (routing spec).} Three-way routing: temporal
recipe only for the $m > 0.4$ core; spatial recipe for $m < -0.2$ (281 heads);
middle $\sim$69 heads stay on the L0 recipe. \textbf{Frozen per-model
\texttt{tau}}, no noise gating, no per-block tables; the source-only
classification pre-pass (Step~1 sequencing) is demoted from v1 default to
fallback. Risk~2 (taxonomy survival) is closed; claim~5 moves to
screening-validated with the CI caveat recorded in the claim itself.

\textbf{Measurement caveats carried forward.} (i) The 47 sampled query rows
are a \emph{fixed panel} (same seeded positions across all captures) ---
positional bias would not average out; extremes are immune, boundary margins
soft. Rows-robustness arm (10\% rows, alternate seeds, one clip) scheduled with
R10. (ii) $n = 5$ videos quantizes per-head flip rates to tenths; the
$< 5\%$ criterion is met at point estimate only. Mitigation: ride-along
profiling in all subsequent runs. (iii) The $m > 0.4$ threshold was chosen on
the data that validates it --- treat as hyperparameter downstream.

Full report with figures: R9 artifact page (margin histogram, 30$\times$12
census grid, stability tables, mechanism probes); raw outputs under
\texttt{five\_bench/r9\_head\_profile/}, \texttt{tau} tensors in
\texttt{evaluation/}.

\end{document}
